\begin{abstract}
India's Central Drugs Standard Control Organisation (CDSCO) publishes monthly lists of drug samples that failed quality tests or were declared spurious. The lists name individual batches, yet a person holding a strip of tablets has no practical way to check them. StripCheck is an offline agent that reads a photo of a medicine pack and decides whether its batch appears on these lists. Two properties of the problem shape the design. Optical character recognition (OCR) misreads single characters with high confidence, so a batch number read from foil cannot be trusted alone. Batch numbers also collide across unrelated products, and sibling batches of one product often differ by a single character, so a matching batch number is weak evidence on its own. StripCheck retrieves candidates by exact and confusion-aware batch search, corroborates each candidate field by field in the manner of record linkage, asks the user to confirm a doubtful character or the batch itself before it raises any alarm, and sends every draft verdict through a rule-based verifier that can abstain. An ingest agent builds the alert index from CDSCO's online portal: \NumAlertRows{} entries from \FirstMonth{} to \LastMonth{}. In these data the median alert appears \LagMedian{} months after manufacture, when \ShelfTwelvePlusPct\% of the flagged batches still have at least a year of shelf life left. \todo{one or two sentences of evaluation results after M7}
\end{abstract}

\begin{IEEEkeywords}
substandard medicines, LLM agents, human-in-the-loop, OCR post-correction, record linkage, abstention
\end{IEEEkeywords}
